% ============================================================================
% Chapter 2: Cardano's Extended UTxO Model and Architecture
% Study Guide for Blockchain Credential Verification
%
% This file contains ONLY chapter content. It must be included from a
% main document that provides the preamble, TikZ libraries, and
% \begin{document} wrapper.
%
% Required TikZ libraries (load in preamble):
%   shapes.geometric, arrows.meta, positioning, fit, calc, backgrounds,
%   decorations.pathreplacing, patterns, chains
%
% Required packages: amsmath, amssymb, xcolor, listings, enumitem, booktabs
% ============================================================================

% --- Chapter-local TikZ styles ---------------------------------------------
\tikzset{
  guidebox/.style={
    draw, rounded corners=4pt, thick, align=center,
    font=\sffamily\small, fill=white, minimum height=1cm
  },
  guidearrow/.style={
    -{Latex[length=2.5mm]}, thick, color=gray!70!black
  },
  guidelabel/.style={
    font=\sffamily\footnotesize, text=gray!50!black
  },
}

% ============================================================================
\chapter{Cardano's Extended UTxO Model and Architecture}
\label{ch:cardano-architecture}
% ============================================================================

In Chapter~1 we established the foundational concepts of blockchain
technology: cryptographic hashing, public-key cryptography, consensus
mechanisms, and the distinction between the account model (Ethereum) and
the UTxO model (Bitcoin). We saw that Bitcoin's Unspent Transaction
Output model treats value like physical coins---each output is a
discrete, indivisible unit that must be consumed entirely when spent,
with any excess returned as ``change.'' This chapter takes that
foundation and builds upon it, exploring how Cardano extends the UTxO
model into something far more powerful: a programmable, data-rich ledger
capable of supporting complex credential verification workflows.

Understanding Cardano's architecture is not merely academic. Every
design decision in the credential verification system described in later
chapters---from how license tokens are minted, to how revocation
propagates, to how transaction costs are predicted---flows directly from
the architectural properties covered here. By the end of this chapter
you will understand \emph{why} Cardano was chosen for this system, not
just \emph{that} it was chosen.


% ============================================================================
\section{The Extended UTxO (eUTxO) Model}
\label{sec:eutxo}
% ============================================================================

\subsection{From Bitcoin's UTxO to Cardano's eUTxO}

Recall from Chapter~1 that Bitcoin's UTxO model works like a collection
of sealed envelopes, each containing a specific amount of currency and a
lock that only the recipient's private key can open. When Alice sends Bob
3~BTC, she does not modify a balance in an account database. Instead, she
\emph{consumes} one or more existing UTxOs she controls and
\emph{produces} new UTxOs: one paying Bob 3~BTC, and possibly another
returning change to herself. The consumed UTxOs cease to exist; the new
UTxOs become available for future spending. This is fundamentally
different from the account model, where a transfer simply decrements
one balance and increments another.

Bitcoin's UTxOs, however, are limited. Each UTxO carries only two pieces
of information: a value (how much BTC it holds) and a locking script
(who can spend it, typically requiring a specific public key signature).
There is no way to attach arbitrary data to a UTxO, no way for a
spending transaction to provide arguments to the locking script, and the
locking script itself has only minimal awareness of the transaction that
attempts to spend it.

Cardano's Extended UTxO model preserves the core advantages of the UTxO
approach---determinism, parallelism, and simple local validation---while
adding three critical extensions that transform UTxOs from simple value
carriers into programmable data containers. These three extensions are
\textbf{Datums}, \textbf{Redeemers}, and \textbf{Script Context}.

\subsection{Datums: Data Attached to UTxOs}

A \textbf{Datum} is an arbitrary piece of structured data that can be
attached to any UTxO. Think of it as a label glued to the outside of the
sealed envelope. In Bitcoin, the envelope only tells you how much money
is inside and whose key opens it. In Cardano, the envelope can also
carry a data payload---a license number, an expiration date, a list of
authorized signers, a document hash, or any other structured information
that a validator script might need to examine.

Datums are encoded using Cardano's Plutus data format (CBOR-serialized
algebraic data types), and they can range from simple integers to deeply
nested records with multiple fields. For credential verification, datums
are the mechanism that turns a native token from a bare proof of
existence into a rich credential record. A license NFT's reference token,
for example, carries an inline datum containing the license type,
jurisdiction, issue date, licensee identity, and a mutable status
field---all readable by any party querying the blockchain.

Before the Vasil hard fork (September 2022), datums had to be provided
separately from the UTxO and matched by hash, which required the
spending party to know the datum's content in advance. The introduction
of \textbf{inline datums} (CIP-32) changed this: the datum is stored
directly within the UTxO on the ledger, making it readable by anyone
without requiring off-chain coordination. This is what makes instant
credential verification possible---a verifier simply reads the inline
datum from the reference token UTxO.

\subsection{Redeemers: Data Provided When Spending}

A \textbf{Redeemer} is a piece of data provided by the party attempting
to spend a script-locked UTxO. If the datum is the label on the
envelope, the redeemer is the message you hand to the guard who decides
whether to let you open it. The redeemer does not exist on the
blockchain until the moment a spending transaction is submitted---it is
transaction-specific, ephemeral data that the validator script receives
as an argument.

Redeemers enable a single validator script to support multiple
operations. For example, the Signature Collection Validator in our
credential system accepts three different redeemer values: ``Collect''
(a signer deposits tokens), ``Finalize'' (the authority completes the
process), and ``Reclaim'' (a signer cancels and retrieves tokens). The
script examines the redeemer to determine which set of validation rules
to apply, then checks the transaction against those rules. This is
analogous to a single contract with multiple clauses---the redeemer
selects which clause is being invoked.

The redeemer pattern is powerful because it separates the \emph{what}
(the action being performed) from the \emph{how} (the validation logic).
A single script address can manage complex multi-step workflows without
requiring separate contracts for each step.

\subsection{Script Context: The Transaction in Full}

The \textbf{Script Context} is the complete set of information about the
transaction that is made available to the validator script during
execution. This includes all inputs being consumed, all outputs being
produced, all minting and burning operations, all signatures present,
the valid time range (slot interval), and the transaction fee. The script
context is what gives Cardano validators their power: they can inspect
the \emph{entire} transaction, not just the single UTxO being spent.

This is a crucial distinction from Bitcoin's Script, which can only see
the specific input being validated. A Cardano validator can enforce rules
like ``this UTxO may only be spent if the transaction also produces an
output paying at least 50 ADA to address X'' or ``this UTxO may only be
consumed if the transaction mints exactly one token under policy Y.''
For credential systems, script context enables validators to verify that
a signer's deposit transaction includes the correct tokens, references
the correct document hash, and falls within the validity window---all in
a single atomic check.

\subsection{Determinism: The Key Advantage}

The combination of datums, redeemers, and script context produces a
property that is perhaps the eUTxO model's most important advantage over
account-based models: \textbf{deterministic execution}. Given the same
inputs, datum, redeemer, and script context, a Cardano validator always
produces the same result, regardless of what other transactions are
happening on the network at the same time.

In Ethereum's account model, the outcome of a smart contract call
depends on the global state at the moment of execution, which can change
between the time a transaction is submitted and when it is included in a
block. This leads to failed transactions that still consume gas,
front-running attacks where miners reorder transactions for profit, and
unpredictable costs. In Cardano's eUTxO model, you can validate a
transaction completely off-chain before submitting it. If it validates
locally, it \emph{will} validate on-chain. The fee is known exactly. The
outcome is guaranteed. For a credential verification system handling
professional licenses and legal signatures, this predictability is not
merely convenient---it is essential.

\begin{figure}[htbp]
\centering
\begin{tikzpicture}[
  scale=0.92,
  every node/.style={transform shape},
]
  % ── Consumed UTxO (input) ──────────────────────────────────────
  \node[guidebox, fill=blue!8, minimum width=5.2cm, minimum height=3.8cm]
    (input) at (-4.2, 0) {};
  \node[font=\sffamily\small\bfseries, text=blue!60!black]
    at (-4.2, 1.55) {Input UTxO (eUTxO)};
  \draw[thick, blue!40] (-6.6, 1.2) -- (-1.8, 1.2);

  % Value area
  \node[guidebox, fill=green!8, minimum width=4.4cm, minimum height=0.7cm]
    at (-4.2, 0.7) {\textbf{Value:} 10 ADA + License NFT};

  % Datum area
  \node[guidebox, fill=yellow!8, minimum width=4.4cm, minimum height=1.0cm]
    at (-4.2, -0.25) {%
      \textbf{Datum (inline)}\\[1pt]
      \footnotesize\texttt{\{type: "PE", status: "active"\}}};

  % Locking script label
  \node[guidebox, fill=purple!8, minimum width=4.4cm, minimum height=0.6cm]
    at (-4.2, -1.15) {\textbf{Locked by:} Validator Script};

  % ── Spending Transaction ───────────────────────────────────────
  \node[guidebox, fill=orange!8, minimum width=4.4cm, minimum height=4.4cm]
    (tx) at (3.0, 0) {};
  \node[font=\sffamily\small\bfseries, text=orange!60!black]
    at (3.0, 1.8) {Spending Transaction};
  \draw[thick, orange!40] (0.95, 1.45) -- (5.05, 1.45);

  % Redeemer
  \node[guidebox, fill=red!6, minimum width=3.6cm, minimum height=0.7cm]
    at (3.0, 0.9) {\textbf{Redeemer:} \texttt{Collect}};

  % Script Context
  \node[guidebox, fill=gray!10, minimum width=3.6cm, minimum height=1.4cm]
    at (3.0, -0.15) {%
      \textbf{Script Context}\\[1pt]
      \footnotesize Inputs, Outputs, Mints\\
      \footnotesize Signatures, Slot Range};

  % Output
  \node[guidebox, fill=green!8, minimum width=3.6cm, minimum height=0.7cm]
    at (3.0, -1.3) {\textbf{Output:} New UTxO(s)};

  % ── Arrows ─────────────────────────────────────────────────────
  \draw[guidearrow, line width=1.2pt]
    (input.east) -- node[above, guidelabel] {consumed by} (tx.west);

  % ── Validator decision ─────────────────────────────────────────
  \node[guidebox, fill=green!15, minimum width=2.8cm, minimum height=0.8cm,
        font=\sffamily\footnotesize\bfseries]
    (accept) at (3.0, -2.8) {Validator: ACCEPT?};

  \draw[guidearrow, dashed] (tx.south) -- (accept.north);

  % Annotations
  \draw[decorate, decoration={brace, amplitude=6pt, raise=3pt},
        thick, blue!50!black]
    (-6.7, 1.9) -- (-6.7, -1.6)
    node[midway, left=10pt, font=\sffamily\scriptsize, text=blue!50!black,
         align=right, text width=2.2cm]
    {Extended UTxO:\\value + datum +\\locking script};

  \draw[decorate, decoration={brace, amplitude=6pt, mirror, raise=3pt},
        thick, orange!50!black]
    (5.3, 1.9) -- (5.3, -1.7)
    node[midway, right=10pt, font=\sffamily\scriptsize, text=orange!50!black,
         align=left, text width=2.4cm]
    {Transaction provides\\redeemer + context\\to the validator};
\end{tikzpicture}
\caption{Anatomy of an Extended UTxO and a spending transaction. The input
UTxO carries a value, an inline datum, and is locked by a validator script.
The spending transaction provides a redeemer (the action to perform) and
the script context (complete transaction data). The validator examines all
three---datum, redeemer, context---to decide whether to accept or reject
the spend.}
\label{fig:eutxo-anatomy}
\end{figure}


% ============================================================================
\section{Native Multi-Asset Framework}
\label{sec:native-assets}
% ============================================================================

\subsection{Tokens Without Smart Contracts}

One of Cardano's most distinctive architectural features---and a primary
reason it was chosen for the credential verification system---is its
\textbf{native multi-asset framework}. On Ethereum, creating a custom
token requires deploying a smart contract (ERC-20 for fungible tokens,
ERC-721 for non-fungible tokens). Every transfer of that token requires
executing the smart contract, consuming gas, and relying on the contract
code being correct. A bug in an ERC-20 contract can freeze or drain
every token it manages, as numerous exploits have demonstrated.

On Cardano, custom tokens---both fungible and non-fungible---exist at
the \emph{ledger level}, the same level as ADA itself. They are tracked
by the same UTxO accounting rules, validated by the same consensus
mechanism, and transferred with the same transaction logic. No smart
contract execution is required to transfer a native token from one
address to another. The token simply moves from one UTxO to another,
just like ADA. This means that native tokens inherit all of the security
guarantees of the base currency: they cannot be double-spent, they
cannot be created out of thin air (except by an authorized minting
policy), and they cannot be frozen by a buggy contract.

For a credential verification system, this architecture is
transformative. A professional license NFT on Cardano is not a record in
a smart contract's internal storage---it is a first-class asset on the
ledger, as ``real'' as ADA itself. Transferring it costs only a standard
transaction fee (\textasciitilde0.17~ADA), not a smart contract
execution fee. Verifying its existence requires only a UTxO query, not a
contract call.

\subsection{Policy ID and Asset Name}

Every native token on Cardano is identified by a two-part identifier:

\begin{description}[leftmargin=1.5cm, labelwidth=1.3cm]
  \item[Policy ID] A 28-byte (56 hex character) Blake2b-224 hash of the
    \textbf{minting policy script}. The minting policy is the script
    that defines who can create (mint) or destroy (burn) tokens under
    this policy. Because the Policy ID is a cryptographic hash of the
    script, it is deterministic and unforgeable---anyone can independently
    verify which script produced a given Policy ID. If the California
    Board of Professional Engineers publishes their Policy ID, any
    verifier can confirm that a license token was minted under that exact
    policy.

  \item[Asset Name] A human-readable (or arbitrary byte string) name
    within a policy, up to 32 bytes. For credential tokens, this
    typically encodes the credential type or identifier, such as
    \texttt{LIC\_PE\_2026\_00042} or \texttt{SIG\_TOKEN}. The
    combination of Policy ID and Asset Name uniquely identifies every
    token on the Cardano ledger.
\end{description}

This two-part identification scheme is analogous to a domain name
system: the Policy ID is like the domain (controlled by the authority),
and the Asset Name is like a path within that domain (chosen by the
authority for each token type). A forger cannot create a token with the
California Board's Policy ID because they do not possess the authority
key that the minting policy requires.

\subsection{Multi-Asset Bundles in UTxOs}

A single Cardano UTxO can carry \emph{multiple} tokens alongside ADA.
This is called a \textbf{multi-asset bundle} (or simply a ``value
bundle''). Unlike Bitcoin, where each UTxO carries only the native
currency, a Cardano UTxO's value field is a nested map:

\begin{quote}
\texttt{Value = Map<PolicyId, Map<AssetName, Quantity>>}
\end{quote}

A single UTxO might contain 5~ADA, plus 100 signature tokens under
Policy~A, plus 1 license NFT under Policy~B, plus 1 validity token
under Policy~C. This bundling capability is essential for the credential
system's signature workflow, where a signer must deposit multiple token
types in a single atomic transaction.

\begin{figure}[htbp]
\centering
\begin{tikzpicture}[
  scale=0.92,
  every node/.style={transform shape},
]
  % ── Main UTxO container ────────────────────────────────────────
  \node[guidebox, fill=blue!8, minimum width=10.2cm, minimum height=6.2cm]
    (utxo) at (0, 0) {};
  \node[font=\sffamily\bfseries, text=blue!60!black]
    at (0, 2.75) {UTxO: Multi-Asset Bundle};
  \draw[thick, blue!40] (-4.9, 2.4) -- (4.9, 2.4);

  % ── Address ────────────────────────────────────────────────────
  \node[guidebox, fill=gray!10, minimum width=9.4cm, minimum height=0.6cm]
    at (0, 1.9)
    {\textbf{Address:}
     \texttt{addr\_test1qz8fg...} (licensee's base address)};

  % ── ADA value ──────────────────────────────────────────────────
  \node[guidebox, fill=green!12, minimum width=9.4cm, minimum height=0.7cm]
    (ada) at (0, 1.05)
    {\textbf{Lovelace (ADA):} 5{,}000{,}000 (5.0 ADA)};

  % ── Token rows ─────────────────────────────────────────────────
  % Token 1: License NFT
  \node[guidebox, fill=orange!8, minimum width=9.4cm, minimum height=1.0cm]
    (tok1) at (0, 0.0) {};
  \node[font=\sffamily\footnotesize, anchor=west] at (-4.3, 0.25)
    {\textbf{Policy:} \texttt{a1b2c3...} (CA Board PE)};
  \node[font=\sffamily\footnotesize, anchor=west] at (-4.3, -0.15)
    {Asset: \texttt{LIC\_PE\_2026\_00042} \quad Qty: \textbf{1} (NFT)};

  % Token 2: Signature tokens
  \node[guidebox, fill=purple!8, minimum width=9.4cm, minimum height=1.0cm]
    (tok2) at (0, -1.15) {};
  \node[font=\sffamily\footnotesize, anchor=west] at (-4.3, -0.9)
    {\textbf{Policy:} \texttt{a1b2c3...} (CA Board PE)};
  \node[font=\sffamily\footnotesize, anchor=west] at (-4.3, -1.3)
    {Asset: \texttt{SIG\_TOKEN} \quad Qty: \textbf{97} (fungible)};

  % Token 3: Validity token
  \node[guidebox, fill=yellow!8, minimum width=9.4cm, minimum height=1.0cm]
    (tok3) at (0, -2.3) {};
  \node[font=\sffamily\footnotesize, anchor=west] at (-4.3, -2.05)
    {\textbf{Policy:} \texttt{d4e5f6...} (Validity Authority)};
  \node[font=\sffamily\footnotesize, anchor=west] at (-4.3, -2.45)
    {Asset: \texttt{VAL\_2026\_Q2} \quad Qty: \textbf{1} (time-bounded)};

  % ── Braces for value map ───────────────────────────────────────
  \draw[decorate, decoration={brace, amplitude=8pt, mirror, raise=4pt},
        thick, gray!60]
    (5.1, 1.4) -- (5.1, -2.8)
    node[midway, right=14pt, font=\sffamily\scriptsize, text=gray!50!black,
         align=left, text width=2.6cm]
    {Multi-asset value:\\[2pt]
     \texttt{Map<PolicyId,}\\
     \texttt{\ Map<AssetName,}\\
     \texttt{\ \ Quantity>{}>}};
\end{tikzpicture}
\caption{A single Cardano UTxO carrying a multi-asset bundle: ADA (required
for the minimum UTxO deposit), a license NFT, 97 fungible signature tokens,
and a time-bounded validity token. All reside in a single UTxO under the
licensee's address. The tokens span two different Policy IDs, demonstrating
that a UTxO can hold tokens from multiple authorities.}
\label{fig:multi-asset-utxo}
\end{figure}


% ============================================================================
\section{Minimum UTxO (minUTxO) Deposit}
\label{sec:min-utxo}
% ============================================================================

\subsection{The Problem: Ledger Bloat}

Every UTxO on the Cardano blockchain consumes storage space on every
full node in the network. Unlike account-based systems where an empty
account can be cheaply represented (or garbage-collected), every UTxO
persists until explicitly consumed by a transaction. If creating UTxOs
were free, an attacker could create billions of tiny, worthless UTxOs to
bloat the ledger, degrading performance for all participants. This is
known as the \textbf{dust attack} problem.

Cardano addresses this with the \textbf{minimum UTxO deposit}
requirement: every UTxO must contain a minimum amount of ADA,
proportional to the size of the UTxO (including any native tokens and
datums it carries). The formula is based on a protocol parameter called
\texttt{coinsPerUTxOByte} (currently 4,310 lovelace per byte). In
practice, this works out to approximately:

\begin{itemize}[nosep]
  \item \textbf{\textasciitilde1.0 ADA} for a simple ADA-only UTxO
  \item \textbf{\textasciitilde1.5 ADA} for a UTxO carrying a single
    native token
  \item \textbf{\textasciitilde2.0 ADA} for a UTxO carrying multiple
    tokens or a datum
  \item \textbf{\textasciitilde2.5--5.0 ADA} for a UTxO with a large
    inline datum (e.g., CIP-68 reference tokens)
\end{itemize}

\subsection{The Apartment Deposit Analogy}

The minUTxO deposit works exactly like a security deposit on an
apartment. When you move in (create a UTxO carrying native tokens), the
landlord (the Cardano protocol) requires you to put down a deposit (the
minimum ADA). This deposit is not a fee---it is not spent or consumed.
It sits in the UTxO, locked alongside your tokens. When you move out
(consume the UTxO by spending its contents), you get the deposit back.
The ADA returns to you as part of the transaction's change output.

This has a direct impact on the economics of the credential system.
Every license NFT, every batch of signature tokens, and every validity
token requires an accompanying ADA deposit of approximately 1.5--2.0
ADA. For a board issuing 10,000 licenses, this means roughly 20,000~ADA
(\textasciitilde\$6,000 at \$0.30/ADA) must be available as working
capital for deposits, even though the ADA is fully recoverable when
tokens are burned or consolidated.

\subsection{Why This Matters for Credential Systems}

Understanding the minUTxO deposit is crucial for three reasons:

\begin{enumerate}[nosep]
  \item \textbf{Cost Planning.} While the deposit is recoverable, it
    represents locked capital. A licensing board must budget for both
    transaction fees (consumed, non-recoverable) and minUTxO deposits
    (locked, recoverable).
  \item \textbf{Token Consolidation.} If a licensee accumulates many
    small UTxOs (e.g., receiving signature tokens in separate batches),
    the total locked ADA can grow substantially. Periodic consolidation
    of UTxOs reduces the total deposit overhead.
  \item \textbf{Burn Recovery.} When a credential is revoked and the
    token burned, the minUTxO deposit is returned. This creates a
    natural economic incentive to clean up expired or revoked
    credentials, keeping the ledger tidy.
\end{enumerate}


% ============================================================================
\section{Slots, Epochs, and Time}
\label{sec:time}
% ============================================================================

\subsection{Cardano's Time Model}

Cardano does not use wall-clock time directly. Instead, time on the
blockchain is measured in \textbf{slots}---discrete, sequential time
intervals. Understanding this time model is essential because the
credential system uses slot-based expiry for time-bounded validity
tokens, and every transaction specifies a valid slot range.

\begin{description}[leftmargin=2.2cm, labelwidth=2cm]
  \item[Slot] The fundamental unit of time on Cardano. Each slot is
    exactly \textbf{1 second} long. Slots are numbered sequentially from
    the genesis block. At the time of writing, the Cardano mainnet is at
    slot \textasciitilde140,000,000+. A slot leader (the stake pool
    selected to produce the next block) has exactly one slot in which to
    create and propagate their block.

  \item[Block] Not every slot produces a block. On average, a block is
    produced every \textbf{\textasciitilde20 seconds} (the active slot
    coefficient is 0.05, meaning 1 in 20 slots produces a block). This
    20-second average block time is a protocol parameter that balances
    throughput against propagation delay across the global network.

  \item[Epoch] An epoch is a period of \textbf{432,000 slots (5 days)}.
    Epochs are the administrative unit of Cardano: stake delegation
    changes, reward distribution, and protocol parameter updates all
    take effect at epoch boundaries. For credential systems, epochs
    provide a natural granularity for periodic operations like dues
    payment deadlines.
\end{description}

\subsection{Slot-Based Transaction Validity}

Every Cardano transaction specifies a \textbf{validity interval}:
a range of slots during which the transaction is valid. A transaction
with \texttt{invalid\_before = 100000} and
\texttt{invalid\_hereafter = 100500} can only be included in a block
produced during slots 100,000 through 100,499. If a block producer
encounters the transaction outside this window, it is rejected.

This mechanism is critical for time-bounded credentials. A validity
token's expiry is expressed as a slot number in its datum. The validator
script checks the transaction's validity interval against the token's
expiry slot: if the transaction's valid slot range extends beyond the
token's expiry, the validator rejects the spend. This ensures that a
validity token cannot be used to sign a document after it has expired,
even if the licensee still holds it in their wallet.

\subsection{Converting Between Slots and Wall-Clock Time}

Converting between slots and human-readable timestamps is
straightforward but requires knowing the chain's genesis parameters. For
mainnet:

\begin{quote}
\texttt{wall\_clock = genesis\_time + (slot\_number * slot\_length)}\\
where \texttt{genesis\_time = 2017-09-23T21:44:51Z}
and \texttt{slot\_length = 1 second}
\end{quote}

In practice, the PyCardano library and the Blockfrost API handle this
conversion automatically. When the credential system mints a validity
token with a 90-day expiry, it calculates the target slot as
\texttt{current\_slot + (90 * 24 * 60 * 60)}.

\begin{figure}[htbp]
\centering
\begin{tikzpicture}[
  scale=0.88,
  every node/.style={transform shape},
]
  % ── Main timeline axis ─────────────────────────────────────────
  \draw[thick, gray!60, -{Latex[length=3mm]}]
    (-6.5, 0) -- (6.5, 0)
    node[right, font=\sffamily\footnotesize, text=gray!50!black] {time};

  % ── Epoch markers ──────────────────────────────────────────────
  \foreach \x/\lab in {-6/Epoch $n$, 0/Epoch $n{+}1$, 6/Epoch $n{+}2$} {
    \draw[thick, blue!60] (\x, -0.25) -- (\x, 0.25);
    \node[font=\sffamily\footnotesize\bfseries, text=blue!60!black,
          above=4pt] at (\x, 0.25) {\lab};
  }

  % 5 days label
  \draw[{Latex[length=2mm]}-{Latex[length=2mm]}, thick, blue!40]
    (-6, -0.55) -- (0, -0.55)
    node[midway, below=2pt, font=\sffamily\scriptsize, text=blue!40!black]
    {432{,}000 slots = 5 days};

  % ── Slot detail (zoomed inset) ─────────────────────────────────
  \node[draw, rounded corners=3pt, thick, dashed, gray!50,
        minimum width=5.4cm, minimum height=2.8cm, fill=white]
    (inset) at (0, -3.0) {};
  \node[font=\sffamily\footnotesize\bfseries, text=gray!50!black]
    at (0, -1.8) {Slot Detail (zoomed)};

  % Slot ticks
  \foreach \i in {0,...,10} {
    \pgfmathsetmacro{\xpos}{-2.2 + \i * 0.44}
    \draw[thin, gray!40] (\xpos, -2.4) -- (\xpos, -2.7);
  }
  \draw[thick, gray!50, -{Latex[length=2mm]}]
    (-2.6, -2.55) -- (2.8, -2.55);

  % 1-second label
  \draw[{Latex[length=1.5mm]}-{Latex[length=1.5mm]}, thick, green!50!black]
    (-2.2, -2.9) -- (-1.76, -2.9)
    node[midway, below=2pt, font=\sffamily\tiny, text=green!50!black]
    {1 sec};

  % Block markers (every ~20 slots)
  \fill[red!60] (-2.2, -2.55) circle (3pt);
  \fill[red!60] (1.46, -2.55) circle (3pt);  % ~slot 8 of 10 shown

  % Block label
  \node[font=\sffamily\tiny, text=red!50!black, above=3pt]
    at (-2.2, -2.4) {Block};
  \node[font=\sffamily\tiny, text=red!50!black, above=3pt]
    at (1.46, -2.4) {Block};

  % ~20s between blocks
  \draw[{Latex[length=1.5mm]}-{Latex[length=1.5mm]}, thick, red!40]
    (-2.2, -3.55) -- (1.46, -3.55)
    node[midway, below=2pt, font=\sffamily\scriptsize, text=red!40!black]
    {\textasciitilde20 seconds between blocks};

  % Arrow from main timeline to inset
  \draw[guidearrow, dashed, gray!50]
    (-1, -0.55) -- (-1, -1.55);

  % ── Validity token timeline ────────────────────────────────────
  \draw[thick, green!40!black, -{Latex[length=2.5mm]}]
    (-6.5, 2.2) -- (6.5, 2.2);
  \node[font=\sffamily\footnotesize, text=green!40!black, above=2pt]
    at (-6.5, 2.2) {};

  % Mint point
  \fill[green!60!black] (-4.5, 2.2) circle (3.5pt);
  \node[font=\sffamily\scriptsize, text=green!50!black, above=6pt]
    at (-4.5, 2.2) {Token minted};

  % Valid range
  \draw[line width=3pt, green!30]
    (-4.5, 2.2) -- (3.5, 2.2);

  % Expiry point
  \fill[red!60] (3.5, 2.2) circle (3.5pt);
  \node[font=\sffamily\scriptsize, text=red!50!black, above=6pt]
    at (3.5, 2.2) {Expiry slot};

  % Expired range
  \draw[line width=3pt, red!15]
    (3.5, 2.2) -- (6.2, 2.2);

  \node[font=\sffamily\scriptsize, text=green!40!black, below=5pt]
    at (-0.5, 2.2) {\textbf{Valid range}};
  \node[font=\sffamily\scriptsize, text=red!40!black, below=5pt]
    at (5.0, 2.2) {\textbf{Expired}};

  % Label
  \node[font=\sffamily\footnotesize\bfseries, text=green!40!black]
    at (0, 3.0) {Validity Token Lifetime};
\end{tikzpicture}
\caption{Cardano's time model. The upper timeline shows a validity token's
lifetime, from minting to slot-based expiry. The middle timeline shows
epochs (5 days each). The inset zooms into individual slots (1 second each)
and blocks (\textasciitilde20-second intervals). Validators check
transaction validity intervals against token expiry slots to enforce
time-bounded credentials.}
\label{fig:time-model}
\end{figure}


% ============================================================================
\section{Transaction Fees}
\label{sec:fees}
% ============================================================================

\subsection{Deterministic Fee Model}

Cardano's fee model is one of its most practically important features
for production systems. Unlike Ethereum, where the cost of a transaction
is determined by a ``gas auction'' (users bid for block space, and
actual costs depend on network congestion at the moment of inclusion),
Cardano uses a \textbf{deterministic, formula-based fee model}:

\[
  \text{fee} = A \times \text{tx\_size (bytes)} + B
\]

where $A$ and $B$ are protocol parameters. At current mainnet values:
\begin{itemize}[nosep]
  \item $A = 44$ lovelace per byte (the ``per-byte coefficient'')
  \item $B = 155{,}381$ lovelace (the ``constant fee''; the minimum fee
    for any transaction)
\end{itemize}

For a typical native script minting transaction of \textasciitilde400
bytes:
\[
  \text{fee} = 44 \times 400 + 155{,}381 = 173{,}181 \text{ lovelace}
  \approx 0.17 \text{ ADA}
\]

This means the \emph{exact fee is known before the transaction is
submitted}. There is no gas estimation, no fee spike due to network
congestion, and no failed transaction that still costs money. You build
the transaction, calculate its serialized size, apply the formula, and
the result is guaranteed. For a credential system handling thousands of
minting and signing operations, this predictability enables precise
budgeting and cost forecasting.

\subsection{Plutus Script Execution Fees}

Transactions that execute Plutus validator scripts incur additional fees
based on two execution resource metrics: \textbf{CPU steps} and
\textbf{memory units}. The Plutus fee formula adds:

\[
  \text{script\_fee} = \text{cpu\_price} \times \text{cpu\_steps}
    + \text{mem\_price} \times \text{mem\_units}
\]

The total transaction fee is the maximum of the size-based fee and the
execution-based fee. For credential system operations:

\begin{itemize}[nosep]
  \item \textbf{Native script operations} (minting, simple transfers):
    \textasciitilde0.17~ADA. No Plutus execution, so only the size-based
    fee applies.
  \item \textbf{Plutus V2 operations} (signature collection, dues
    enforcement, CIP-68 datum updates): \textasciitilde0.3--0.5~ADA.
    The script execution adds CPU and memory costs.
  \item \textbf{Complex Plutus operations} (multi-input finalization
    consuming 5+ UTxOs): \textasciitilde0.5--1.5~ADA. More inputs mean
    more script context for the validator to process.
\end{itemize}

\subsection{Comparison with Ethereum Gas}

The practical difference is dramatic. During high network activity on
Ethereum (e.g., a popular NFT mint), gas prices can spike from a base of
\textasciitilde20~gwei to 500+~gwei, making a \$2 mint cost \$50 or
more. A transaction that fails validation \emph{still consumes gas} on
Ethereum---the user pays for the failed attempt. On Cardano, fee spikes
do not occur because fees depend only on transaction size and script
execution cost, not on network congestion. And because transactions are
validated locally before submission, a malformed transaction never
reaches the chain and costs nothing.

For a licensing board processing hundreds of minting transactions per
day, this difference is not marginal---it is the difference between
predictable operating costs and unpredictable, potentially order-of-magnitude cost variability.


% ============================================================================
\section{Addresses}
\label{sec:addresses}
% ============================================================================

\subsection{Address Types}

A Cardano address is not simply a hash of a public key (as in Bitcoin).
It is a structured data object that encodes multiple pieces of
information. Understanding address construction is important because the
credential system uses address properties to determine wallet roles and
enforce access control.

Cardano defines several address types, three of which are relevant to
credential systems:

\begin{description}[leftmargin=2.8cm, labelwidth=2.6cm]
  \item[Base Address] The most common address type. It contains two
    credential hashes: a \textbf{payment key hash (PKH)} and a
    \textbf{stake key hash (SKH)}. The PKH controls spending (who can
    consume UTxOs at this address), and the SKH controls stake
    delegation (which stake pool receives delegation from this address's
    ADA). Base addresses are used for all licensee and signer wallets in
    the credential system, as they support both transactions and staking
    rewards.

  \item[Enterprise Address] Contains only a payment key hash---no stake
    component. Enterprise addresses cannot delegate stake or receive
    staking rewards. They are useful for script addresses and
    application-specific wallets where staking is irrelevant. The
    credential system's validator script addresses are enterprise
    addresses.

  \item[Pointer Address] Contains a payment key hash and a ``pointer''
    to a stake key registration on the chain (identified by slot, transaction
    index, and certificate index). Pointer addresses are more compact
    than base addresses but are rarely used in practice and are not
    employed in the credential system.
\end{description}

\subsection{Key Hashes: PKH and SKH}

The payment key hash (PKH) and stake key hash (SKH) are both
28-byte Blake2b-224 hashes of the respective verification keys. These
hashes serve as the fundamental identity unit on Cardano:

\begin{itemize}[nosep]
  \item The \textbf{PKH} identifies \emph{who can spend} UTxOs at an
    address. When a validator script checks whether the correct signer
    is depositing tokens, it compares the transaction's signatories
    against the expected PKH in the datum.
  \item The \textbf{SKH} identifies \emph{who controls delegation}. It
    is relevant for staking rewards but does not affect transaction
    authorization.
\end{itemize}

In the credential system, the authority's PKH is embedded in minting
policies (``only the holder of this PKH can mint''), signer PKHs are
recorded in work product datums (``these specific PKHs must deposit
tokens''), and the licensee's PKH is recorded in license metadata for
identity binding.

\subsection{Bech32 Encoding}

Cardano addresses are displayed in \textbf{Bech32} encoding, a
human-readable format that includes an error-detecting checksum. Mainnet
addresses start with \texttt{addr1...} and testnet addresses with
\texttt{addr\_test1...}. The Bech32 encoding ensures that a mistyped
character in an address is detected before a transaction is submitted,
preventing accidental loss of funds. The prefix also prevents
accidentally sending mainnet transactions on the testnet or vice versa.

\begin{figure}[htbp]
\centering
\begin{tikzpicture}[
  scale=0.88,
  every node/.style={transform shape},
]
  % ── Mnemonic at top ────────────────────────────────────────────
  \node[guidebox, fill=gray!10, minimum width=8cm, minimum height=0.9cm]
    (mnemonic) at (0, 5.2)
    {\textbf{24-word BIP-39 Mnemonic} (entropy source)};

  % ── HD Derivation ──────────────────────────────────────────────
  \node[guidebox, fill=blue!8, minimum width=3.4cm, minimum height=1.0cm]
    (paypath) at (-2.5, 3.6)
    {\textbf{Payment Key}\\
     \footnotesize\texttt{m/1852'/1815'/0'/0/0}};
  \node[guidebox, fill=blue!8, minimum width=3.4cm, minimum height=1.0cm]
    (stakepath) at (2.5, 3.6)
    {\textbf{Stake Key}\\
     \footnotesize\texttt{m/1852'/1815'/0'/2/0}};

  \draw[guidearrow] (mnemonic.south west) ++(0.8,0) -- (paypath.north);
  \draw[guidearrow] (mnemonic.south east) ++(-0.8,0) -- (stakepath.north);

  % ── Verification Keys ──────────────────────────────────────────
  \node[guidebox, fill=green!8, minimum width=3.4cm, minimum height=0.8cm]
    (payvk) at (-2.5, 2.1)
    {\textbf{Payment VK} (Ed25519, 32\,B)};
  \node[guidebox, fill=green!8, minimum width=3.4cm, minimum height=0.8cm]
    (stakevk) at (2.5, 2.1)
    {\textbf{Stake VK} (Ed25519, 32\,B)};

  \draw[guidearrow] (paypath.south) -- (payvk.north);
  \draw[guidearrow] (stakepath.south) -- (stakevk.north);

  % ── Hashing ────────────────────────────────────────────────────
  \node[guidebox, fill=orange!8, minimum width=3.4cm, minimum height=0.8cm]
    (pkh) at (-2.5, 0.7)
    {\textbf{PKH} (Blake2b-224, 28\,B)};
  \node[guidebox, fill=orange!8, minimum width=3.4cm, minimum height=0.8cm]
    (skh) at (2.5, 0.7)
    {\textbf{SKH} (Blake2b-224, 28\,B)};

  \draw[guidearrow] (payvk.south) --
    node[right=2pt, guidelabel] {hash} (pkh.north);
  \draw[guidearrow] (stakevk.south) --
    node[right=2pt, guidelabel] {hash} (skh.north);

  % ── Address Assembly ───────────────────────────────────────────
  \node[guidebox, fill=purple!8, minimum width=8.6cm, minimum height=1.4cm]
    (addr) at (0, -0.8) {};
  \node[font=\sffamily\small\bfseries, text=purple!60!black]
    at (0, -0.35) {Base Address Construction};
  \node[font=\sffamily\footnotesize] at (0, -0.85)
    {Header (1\,B) $\|$ PKH (28\,B) $\|$ SKH (28\,B)
     $\|$ Network Tag};

  \draw[guidearrow] (pkh.south) -- (addr.north -| pkh);
  \draw[guidearrow] (skh.south) -- (addr.north -| skh);

  % ── Bech32 Encoding ───────────────────────────────────────────
  \node[guidebox, fill=yellow!8, minimum width=8.6cm, minimum height=0.9cm]
    (bech32) at (0, -2.4)
    {\textbf{Bech32:} \texttt{addr1qy...} (mainnet) /
     \texttt{addr\_test1qz...} (testnet)};

  \draw[guidearrow] (addr.south) --
    node[right=2pt, guidelabel] {Bech32 encode} (bech32.north);

  % ── Enterprise address (side note) ────────────────────────────
  \node[guidebox, fill=red!6, minimum width=3.8cm, minimum height=1.0cm,
        dashed, font=\sffamily\scriptsize]
    (enterprise) at (6.5, -0.8)
    {\textbf{Enterprise Address}\\PKH only, no SKH\\(for script addresses)};

  \draw[guidearrow, dashed, gray!40]
    (addr.east) -- (enterprise.west);

  % ── Legend ─────────────────────────────────────────────────────
  \node[font=\sffamily\tiny, text=gray!40!black, anchor=north west]
    at (-4.5, -3.2)
    {PKH = Payment Key Hash \quad SKH = Stake Key Hash \quad
     VK = Verification Key \quad B = Bytes};
\end{tikzpicture}
\caption{Cardano address construction from mnemonic to Bech32-encoded
address. A 24-word mnemonic seeds HD key derivation (CIP-1852), producing
payment and stake verification keys. Each is hashed with Blake2b-224 to
produce 28-byte key hashes. A base address combines both hashes with a
network tag; an enterprise address uses only the payment key hash.}
\label{fig:address-construction}
\end{figure}


% ============================================================================
\section{UTxO Selection Strategies}
\label{sec:utxo-selection}
% ============================================================================

\subsection{The UTxO Selection Problem}

When building a transaction, you must select which UTxOs to consume as
inputs. This is a non-trivial optimization problem. Consider a licensee
who has received signature tokens in 20 separate transactions, each
creating a separate UTxO. To sign a document (requiring 1 signature
token and 1 validity token), the transaction builder must choose which
UTxO(s) to consume. The choice affects:

\begin{enumerate}[nosep]
  \item \textbf{Transaction size.} Each input adds \textasciitilde40--60
    bytes to the transaction. More inputs mean a larger transaction and
    a higher fee.
  \item \textbf{Change outputs.} If a UTxO contains more value than
    needed, the excess must be returned as a change output, which itself
    requires a minUTxO deposit.
  \item \textbf{UTxO set health.} Poor selection strategies can
    fragment the UTxO set into many tiny outputs, each locking a minUTxO
    deposit, gradually eroding the wallet's effective balance.
  \item \textbf{Transaction size limit.} Cardano enforces a maximum
    transaction size of \textbf{16,384 bytes (16~KB)}. A transaction
    that consumes too many inputs may exceed this limit and fail to
    build.
\end{enumerate}

\subsection{LargestFirstSelector}

The \texttt{LargestFirstSelector} strategy sorts all available UTxOs by
value (largest first) and greedily selects UTxOs until the required
value is covered. This is the default strategy in PyCardano and the one
used by the credential system.

\textbf{Advantages:}
\begin{itemize}[nosep]
  \item Minimizes the number of inputs (and therefore transaction size)
    by preferring large UTxOs.
  \item Produces fewer, larger change outputs, reducing UTxO
    fragmentation over time.
  \item Simple and predictable---the same inputs produce the same
    selection.
\end{itemize}

\textbf{Disadvantages:}
\begin{itemize}[nosep]
  \item Large UTxOs are consumed first, potentially leaving a ``long
    tail'' of small UTxOs that are individually too small to cover
    transaction fees.
  \item Deterministic selection can lead to repeated use of the same
    UTxOs, causing temporary contention in high-throughput scenarios.
\end{itemize}

\subsection{RandomImproveSelector}

The \texttt{RandomImproveSelector} takes a different approach. It first
makes a random initial selection that covers the required value, then
iteratively attempts to improve the selection by swapping UTxOs to
minimize the change output. The goal is to produce a change output as
close to zero as possible, reducing the number of dust UTxOs over time.

\textbf{Advantages:}
\begin{itemize}[nosep]
  \item Distributes UTxO consumption more evenly across the wallet's
    UTxO set, reducing fragmentation over time.
  \item The randomness reduces contention in scenarios where multiple
    transactions are built concurrently (relevant for high-volume
    minting operations).
  \item Tends to produce a healthier UTxO distribution in the long run.
\end{itemize}

\textbf{Disadvantages:}
\begin{itemize}[nosep]
  \item Non-deterministic: building the same transaction twice may
    select different inputs, complicating testing and debugging.
  \item May select more inputs than \texttt{LargestFirstSelector} for a
    given transaction, increasing fees slightly.
\end{itemize}

\subsection{Practical Implications}

For the credential system, UTxO selection strategy matters most during
two operations:

\begin{enumerate}[nosep]
  \item \textbf{Batch minting.} When minting 100 signature tokens for a
    new licensee, the authority wallet must have a UTxO with enough ADA
    to cover fees and the minUTxO deposit. If the authority's ADA is
    fragmented across dozens of small UTxOs, the transaction builder may
    need to consume many inputs, potentially exceeding the 16~KB limit.
  \item \textbf{Token consolidation.} A licensee who has signed many
    documents will have accumulated change UTxOs. Periodic consolidation
    (sending all UTxOs back to oneself in a single transaction) reduces
    the UTxO count and recovers locked minUTxO deposits.
\end{enumerate}

The credential system uses \texttt{LargestFirstSelector} by default
with a 200,000~lovelace fee buffer, switching to
\texttt{RandomImproveSelector} for high-volume operations where
concurrent transaction building might cause contention on the same
large UTxOs.


% ============================================================================
\section{Chapter Summary}
\label{sec:ch2-summary}
% ============================================================================

This chapter has covered the seven architectural pillars of Cardano that
underpin the credential verification system:

\begin{enumerate}[nosep]
  \item \textbf{The eUTxO Model} extends Bitcoin's UTxO with datums
    (data on UTxOs), redeemers (spending arguments), and script context
    (full transaction visibility), enabling deterministic smart contract
    execution.
  \item \textbf{Native Multi-Assets} allow tokens to exist at the
    ledger level without smart contracts, providing the same security
    guarantees as ADA at lower cost.
  \item \textbf{The minUTxO Deposit} prevents ledger bloat by requiring
    a recoverable ADA deposit proportional to UTxO size---a cost that
    must be budgeted but is not consumed.
  \item \textbf{Slots, Epochs, and Time} provide the discrete time
    model that enables slot-based token expiry and time-bounded
    transaction validity.
  \item \textbf{Deterministic Fees} ensure that transaction costs are
    known exactly before submission, enabling precise cost forecasting
    for credential operations.
  \item \textbf{Address Construction} from HD-derived key hashes
    provides the identity primitives (PKH, SKH) used throughout the
    system for access control and credential binding.
  \item \textbf{UTxO Selection Strategies} balance transaction
    efficiency against UTxO set health, with implications for
    high-volume minting and long-term wallet management.
\end{enumerate}

With this architectural foundation in place, Chapter~3 will introduce
the specific token types, minting policies, and smart contracts that
implement the credential verification system on top of these primitives.
